% !TEX root = task-3-8-4.tex
% Editable task bank. Append new BankTask blocks at the end of the document.
\documentclass[12pt,a4paper]{article}
\usepackage{physics-textbook}
\usepackage{graphicx,needspace}
\usetikzlibrary{decorations.pathmorphing}
\newif\ifShowSolutions
\ifdefined\TutorEdition
  \ShowSolutionstrue
\else
  \ShowSolutionsfalse % Change to \ShowSolutionstrue to display solutions.
\fi
\newcounter{banktask}
\newsavebox{\BankTaskBox}
\newcommand{\BankTaskContent}[3]{%
  \noindent\textbf{\thebanktask.}\quad #2\par
  \ifShowSolutions
    \Needspace{50mm}% Keep the solution heading with its opening text or diagram.
    #3\par
    \ifdefined\TutorEdition
      \ifcsname TutorSource#1\endcsname
        \nopagebreak\smallskip{\small\csname TutorSource#1\endcsname\par}%
      \fi
    \fi
  \fi
}
\newcommand{\BankTask}[3]{%
  \refstepcounter{banktask}\par\medskip
  \sbox{\BankTaskBox}{\begin{minipage}{\linewidth}%
    \BankTaskContent{#1}{#2}{#3}\end{minipage}}%
  \ifdim\dimexpr\ht\BankTaskBox+\dp\BankTaskBox\relax>\dimexpr\textheight-20mm\relax
    \Needspace{35mm}\BankTaskContent{#1}{#2}{#3}%
  \else
    \Needspace{\dimexpr\ht\BankTaskBox+\dp\BankTaskBox\relax}%
    \noindent\usebox{\BankTaskBox}\par
  \fi
  \medskip
}
\begin{document}
\BookSubsection{3.8.4}{Скорость направленного движения зарядов}


% Source sdamgia-6931: https://phys-ege.sdamgia.ru/problem?id=6931
\BankTask{sdamgia-6931}{
Пучок медленных электронов массой $m$ с зарядом $e$ разгоняется в электронно-лучевой трубке, проходя большую ускоряющую разность потенциалов $U$. Концентрация электронов в пучке после ускорения равна $n$, площадь поперечного сечения пучка $S$. Установите соответствие между физическими величинами и формулами, по которым их можно определить.\par

\begin{center}\begin{minipage}[t]{.47\linewidth}\textbf{Физическая Величина}\par\smallskip А) скорость электронов в пучке после ускорения\par
Б) сила тока в пучке после ускорения\end{minipage}\hfill\begin{minipage}[t]{.47\linewidth}\textbf{Формула}\par\smallskip $1$) $\sqrt{\frac{ eU}{ 2m }}$\par
$2$) $\sqrt{\frac{ 2eU}{ m }}$\par
$3$) $enS \sqrt{\frac{ 2eU}{ m }}$\par
$4$) $enS \sqrt{\frac{ eU}{ 2m }}$\end{minipage}\end{center}
\par
\begin{center}\small\setlength{\tabcolsep}{3pt}\begin{tabular}{|p{0.45\linewidth}|p{0.45\linewidth}|}\hline
А & Б\\\hline
 & \\\hline
\end{tabular}\end{center}\par
}{
\smallskip\noindent\textbf{Решение}\par
При ускорении в электрическом поле электроны приобретают энергию $eU.$ Таким образом, приравнивая кинетическую энергию электронов и энергию, полученную от поля можем найти скорость:\par
\[\frac{ m v ^2 }{ 2 } =eU.\]
$v = \sqrt{\frac{ 2eU}{ m }}$ — получили формулу под номером $2$.\par
Сила тока может быть вычислена по формуле: $I= \frac{ \Delta q}{ \Delta t } .$ Учтем, что $\Delta q=eN=eS\Delta ln,$ где $N$ — число частиц, $\Delta l$ — малое перемещение электронов за малый интервал времени.\par
Тогда: $I= \frac{ eSn\Delta l}{ \Delta t } =enS v =enS \sqrt{\frac{ 2eU}{ m }}$ — получили формулу под номером $3$.\par
Ответ: $23$.\par
}

\end{document}
